\ifdefined\gkver\else\def\gkver{student}\fi
\documentclass[\gkver]{gaokaozhenti}
\begin{document}

\chapter{2007年重庆理综}
%% paperType: 理综物理部分

\begin{enumerate}
\item
%% number: 14
%% paperName: 2007年重庆理综第 14 题 6 分
%% typeId: 1
%% type: 单选题
%% chapter: 原子和原子核
%% point: 玻尔模型
%% method: 无
%% score: 6
%% degree: 650
%% duplicateId: 0
%% body:
可见光光子的能量在 $1.61\UeV\sim3.10\UeV$ 范围内。若氢原子从高能级跃迁到量子数为 $n$ 的低能级的谱线中有可见光，根据氢原子能级图可判断 $n$ 为\xzanswer{B}

\onepicture[w=4.5cm]{figs/14.png}

\fourchoices[answer=B]
{1}
{2}
{3}
{4}

%% answer: B
%% memo:
\memoanswer{
本题原卷及材料中未提供详解，待补充。
}

\item
%% number: 15
%% paperName: 2007年重庆理综第 15 题 6 分
%% typeId: 1
%% type: 单选题
%% chapter: 电路
%% point: 非纯电阻电路
%% method: 无
%% score: 6
%% degree: 450
%% duplicateId: 0
%% body:
汽车电动机启动时车灯会瞬时变暗，如图，在打开车灯的情况下，电动机未启动时电流表读数为 $10\UA$，电动机启动时电流表读数为 $58\UA$，若电源电动势为 $12.5\UV$，内阻为 $0.05\UO$，电流表内阻不计，则因电动机启动，车灯的电功率降低了\xzanswer{B}

\onepicture[w=5.5cm]{figs/15.png}

\fourchoices[answer=B]
{$35.8\UW$}
{$43.2\UW$}
{$48.2\UW$}
{$76.8\UW$}

%% answer: B
%% memo:
\memoanswer{
本题原卷及材料中未提供详解，待补充。
}

\item
%% number: 16
%% paperName: 2007年重庆理综第 16 题 6 分
%% typeId: 1
%% type: 单选题
%% chapter: 电场
%% point: 带电粒子的平衡
%% method: 无
%% score: 6
%% degree: 650
%% duplicateId: 0
%% body:
如图所示，悬挂在 $O$ 点的一根不可伸长的绝缘细线下端有一个带电量不变的小球 A。在两次实验中，均缓慢移动另一带同种电荷的小球 B。当 B 到达悬点 O 的正下方并与 A 在同一水平线上，A 处于受力平衡时，悬线偏离竖直方向的角度为 $\theta$，若两次实验中 B 的电量分别为 $q_{1}$ 和 $q_{2}$，$\theta$ 分别为 $30^{\circ}$ 和 $45^{\circ}$。则 $q_{2}/q_{1}$ 为\xzanswer{C}

\onepicture[w=2.6cm]{figs/16.png}

\fourchoices[answer=C]
{2}
{3}
{$2\sqrt{3}$}
{$3\sqrt{3}$}

%% answer: C
%% memo:
\memoanswer{
本题原卷及材料中未提供详解，待补充。
}

\item
%% number: 17
%% paperName: 2007年重庆理综第 17 题 6 分
%% typeId: 1
%% type: 单选题
%% chapter: 气体性质
%% point: 气体的状态参量
%% method: 无
%% score: 6
%% degree: 450
%% duplicateId: 0
%% body:
为估算池中睡莲叶面承受雨滴撞击产生的平均压强，小明在雨天将一圆柱形水杯置于露台，测得 1 小时内杯中水上升了 $45\Umm$。查询得知，当时雨滴竖直下落速度约为 $12\Ums$。据此估算该压强约为（设雨滴撞击睡莲后无反弹，不计雨滴重力，雨水的密度为 $1\times10^{3}\Ukgmc$）\xzanswer{A}

\fourchoices[answer=A]
{$0.15\UPa$}
{$0.54\UPa$}
{$1.5\UPa$}
{$5.4\UPa$}

%% answer: A
%% memo:
\memoanswer{
本题原卷及材料中未提供详解，待补充。
}

\item
%% number: 18
%% paperName: 2007年重庆理综第 18 题 6 分
%% typeId: 1
%% type: 单选题
%% chapter: 光学
%% point: 光电效应
%% method: 无
%% score: 6
%% degree: 250
%% duplicateId: 0
%% body:
真空中有一平行板电容器，两极板分别由铂和钾（其极限波长分别为 $\lambda_{1}$ 和 $\lambda_{2}$）制成，板面积为 $S$，间距为 $d$。现用波长为 $\lambda$（$\lambda_{1}<\lambda<\lambda_{2}$）的单色光持续照射两板内表面，则电容器的最终带电量 $Q$ 正比于\xzanswer{D}

\fourchoices[answer=D]
{$\frac{d(\lambda-\lambda_{1})}{S\lambda\lambda_{1}}$}
{$\frac{d(\lambda_{2}-\lambda)}{S\lambda\lambda_{2}}$}
{$\frac{S(\lambda-\lambda_{1})}{d\lambda\lambda_{1}}$}
{$\frac{S(\lambda_{2}-\lambda)}{d\lambda\lambda_{2}}$}

%% answer: D
%% memo:
\memoanswer{
本题原卷及材料中未提供详解，待补充。
}

\item
%% number: 19
%% paperName: 2007年重庆理综第 19 题 6 分
%% typeId: 2
%% type: 多选题
%% chapter: 曲线运动
%% point: 天体运动
%% method: 无
%% score: 6
%% degree: 650
%% duplicateId: 0
%% body:
土卫十和土卫十一是土星的两颗卫星，都沿近似为圆周的轨道绕土星运动，其参数如表：

\begin{center}
\begin{tabular}{|c|c|c|c|}
\hline
 & 卫星半径（$\Um$） & 卫星质量（$\Ukg$） & 轨道半径（$\Um$） \\
\hline
土卫十 & $8.90\times10^{4}$ & $2.01\times10^{18}$ & $1.51\times10^{8}$ \\
\hline
土卫十一 & $5.70\times10^{4}$ & $5.60\times10^{17}$ & $1.51\times10^{8}$ \\
\hline
\end{tabular}
\end{center}

两卫星相比，土卫十\xzanswer{AD}

\fourchoices[answer=AD]
{受土星的万有引力较大}
{绕土星做圆周运动的周期较大}
{绕土星做圆周运动的向心加速度较大}
{动能较大}

%% answer: AD
%% memo:
\memoanswer{
本题原卷及材料中未提供详解，待补充。
}

\item
%% number: 20
%% paperName: 2007年重庆理综第 20 题 6 分
%% typeId: 2
%% type: 多选题
%% chapter: 磁场
%% point: 带电粒子在磁场中的运动
%% method: 无
%% score: 6
%% degree: 650
%% duplicateId: 0
%% body:
下列说法正确的是\xzanswer{CD}

\fourchoices[answer=CD]
{正弦交变电流的有效值是最大值的 $\sqrt{2}$ 倍}
{声波是纵波，声源振动越快，声波传播也越快}
{在某介质中，红光折射率比其他色光的小，故红光传播速度比其他色光的大}
{质子和 $\alpha$ 粒子以相同速度垂直进入同一匀强磁场，质子做圆周运动的半径较小}

%% answer: CD
%% memo:
\memoanswer{
本题原卷及材料中未提供详解，待补充。
}

\item
%% number: 21
%% paperName: 2007年重庆理综第 21 题 6 分
%% typeId: 2
%% type: 多选题
%% chapter: 气体性质
%% point: 热力学第一定律
%% method: 无
%% score: 6
%% degree: 450
%% duplicateId: 0
%% body:
氧气钢瓶充气后压强高于外界大气压，假设缓慢漏气时瓶内外温度始终相等且保持不变，忽略氧气分子之间的相互作用。在该漏气过程中瓶内氧气\xzanswer{BC}

\fourchoices[answer=BC]
{分子总数减少，分子总动能不变}
{密度降低，分子平均动能不变}
{吸收热量，膨胀做功}
{压强降低，不对外做功}

%% answer: BC
%% memo:
\memoanswer{
本题原卷及材料中未提供详解，待补充。
}

\item
%% number: 22
%% paperName: 2007年重庆理综第 22 题 11 分
%% typeId: 4
%% type: 实验题
%% chapter: 力物体的平衡
%% point: 胡克定律
%% method: 无
%% score: 11
%% degree: 450
%% duplicateId: 0
%% body:
\begin{enumerate}
	\item
	在“描绘小灯泡的伏安特性曲线”实验中，用导线 a、b、c、d、e、f、g 和 h 按如图所示方式连接电路，电路中所有元器件都完好，且电压表和电流表已调零。闭合开关后：

	\onepicture[w=6.0cm, align=right]{figs/22.png}

	\begin{steps}
		\item
		若电压表的示数为 $2\UV$，电流表的示数为零，小灯泡不亮，则断路的导线为 \tkanswer[1.5cm]{d}；
		\item
		若电压表的示数为零，电流表的示数为 $0.3\UA$，小灯泡亮，则断路的导线为 \tkanswer[1.5cm]{h}；
		\item
		若反复调节滑动变阻器，小灯泡亮度发生变化，但电压表、电流表的示数不能调为零，则断路的导线为 \tkanswer[2cm]{g}。
	\end{steps}

	\item
	建造重庆长江大桥复线桥，将长百米、重千余吨的钢梁从江水中吊起（如图\ref{2007重庆理综22a}所示），施工时采用了将钢梁与水面成一定倾角出水的起吊方案。为了探究该方案的合理性，某研究性学习小组做了两个模拟实验：研究将钢板从水下水平拉出（实验 1）和以一定倾角拉出（实验 2）的过程中总拉力的变化情况。

	\begin{steps}
		\item
		必要的实验器材有：钢板、细绳、水盆、水、支架、刻度尺、计时器和 \tkanswer[2.2cm]{测力计}。
		\item
		根据实验曲线（如图\ref{2007重庆理综22b}），实验 2 中的最大总拉力比实验 1 中的最大总拉力降低了 \tkanswer[2.5cm]{13.3}。
		\item
		根据分子动理论，实验 1 中最大总拉力明显增大的原因是 \tkanswer[3cm]{分子之间存在引力，钢板与水面的接触面积大}。
		\item
		可能导致测量拉力的实验误差的原因有：读数不准、钢板有油污、\tkanswer[3cm]{快速拉出、变速拉出}等等（答出两个即可）。
	\end{steps}
\end{enumerate}

\twopicture[labelprefix={图}, numA=1, numB=2, labelA=2007重庆理综22a, labelB=2007重庆理综22b, wA=4.0cm, wB=3.2cm, align=right]
{figs/22a.png}
{figs/22b.png}

%% answer:
\jdanswer{
\begin{enumerate}
	\item
	① d；② h；③ g。

	\item
	① 测力计（弹簧测力计、力传感器等等）；② 13.3（允许误差 $\pm0.5$）；0.27$\UN$（允许误差 $\pm0.03$）；③ 分子之间存在引力，钢板与水面的接触面积大；④ 快速拉出、变速拉出、出水过程中角度变化、水中有油污、水面波动等等。
\end{enumerate}
}

%% memo:
\memoanswer{
本题原卷及材料中未提供详解，待补充。
}

\item
%% number: 23
%% paperName: 2007年重庆理综第 23 题 16 分
%% typeId: 6
%% type: 计算题
%% chapter: 电磁感应
%% point: 电磁感应受力分析
%% method: 无
%% score: 16
%% degree: 650
%% duplicateId: 0
%% body:
$t=0$ 时，磁场在 $xOy$ 平面内的分布如图所示。其磁感应强度的大小均为 $B_{0}$，方向垂直于 $xOy$ 平面，相邻磁场区域的磁场方向相反。每个同向磁场区域的宽度均为 $l_{0}$。整个磁场以速度 $v$ 沿 $x$ 轴正方向匀速运动。

（1）若在磁场所在区间，$xOy$ 平面内放置一由 $n$ 匝线圈串联而成的矩形导线框 abcd，线框的 bc 边平行于 $x$ 轴，$bc=l_{B}$、$ab=L$，总电阻为 $R$，线框始终保持静止。求：
\begin{enumerate}
	\item
	线框中产生的总电动势大小和导线中的电流大小；
	\item
	线框所受安培力的大小和方向。
\end{enumerate}

（2）该运动的磁场可视为沿 $x$ 轴传播的波，设垂直于纸面向外的磁场方向为正，画出 $t=0$ 时磁感应强度的波形图，并求波长 $\lambda$ 和频率 $f$。

\onepicture[w=7.5cm, align=right]{figs/23.png}

%% answer:
\jdanswer{
\begin{enumerate}
	\item
	① $E=2nB_{0}Lv$，$I=\frac{2nB_{0}Lv}{R}$；② $F=2nB_{0}LI=\frac{4n^{2}B_{0}^{2}L^{2}v}{R}$，方向始终沿 $x$ 轴正方向。

	\item
	$\lambda=2l_{0}$，$f=\frac{v}{2l_{0}}$。
\end{enumerate}
}

%% memo:
\memoanswer{
（1）① 切割磁感线的速度为 $v$，任意时刻线框中电动势大小 $E=2nB_{0}Lv$，导线中的电流大小 $I=\frac{2nB_{0}Lv}{R}$。

② 线框所受安培力的大小为 $F=2nB_{0}LI=\frac{4n^{2}B_{0}^{2}L^{2}v}{R}$，由左手定则判断，线框所受安培力的方向始终沿 $x$ 轴正方向。

（2）磁感应强度的波长和频率分别为 $\lambda=2l_{0}$，$f=\frac{v}{2l_{0}}$。$t=0$ 时磁感应强度的波形图如图。

\onepicture[w=4.0cm, align=right]{figs/23_an.png}
}

\item
%% number: 24
%% paperName: 2007年重庆理综第 24 题 19 分
%% typeId: 6
%% type: 计算题
%% chapter: 电场
%% point: 带电粒子的直线运动
%% method: 无
%% score: 19
%% degree: 450
%% duplicateId: 0
%% body:
飞行时间质谱仪可通过测量离子飞行时间得到离子的荷质比 $q/m$，如图\ref{2007重庆理综24a}。带正电的离子经电压为 $U$ 的电场加速后进入长度为 $L$ 的真空管 AB，可测得离子飞越 AB 所用时间 $t_{1}$。改进以上方法，如图\ref{2007重庆理综24b}，让离子飞越 AB 后进入场强为 $E$（方向如图）的匀强电场区域 BC，在电场的作用下离子返回 B 端，此时，测得离子从 A 出发后飞行的总时间 $t_{2}$（不计离子重力）。

\begin{enumerate}
	\item
	忽略离子源中离子的初速度，① 用 $t_{1}$ 计算荷质比；② 用 $t_{2}$ 计算荷质比。
	\item
	离子源中相同荷质比离子的初速度不尽相同，设两个荷质比都为 $\frac{q}{m}$ 的离子在 A 端的速度分别为 $v$ 和 $v'$（$v\neq v'$），在改进后的方法中，它们飞行的总时间通常不同，存在时间差 $\Delta t$，可通过调节电场 $E$ 使 $\Delta t=0$。求此时 $E$ 的大小。
\end{enumerate}

\twopicture[labelA=2007重庆理综24a, labelB=2007重庆理综24b, wA=5.0cm, wB=5.5cm, align=right]
{figs/24a.png}
{figs/24b.png}

%% answer:
\jdanswer{
\begin{enumerate}
	\item
	① $\frac{L^{2}}{2Ut_{1}^{2}}$；② $\frac{1}{2U}\left(L+\frac{4U}{E}\right)^{2}\frac{1}{t_{2}^{2}}$。

	\item
	$E=\frac{2mvv'}{qL}$。
\end{enumerate}
}

%% memo:
\memoanswer{
（1）① 设离子带电量为 $q$，质量为 $m$，经电场加速后的速度为 $v$，则 $qU=\frac{1}{2}mv^{2}$。离子飞越真空管，在 AB 做匀速直线运动，则 $L=vt_{1}$，解得离子荷质比 $\frac{q}{m}=\frac{L^{2}}{2Ut_{1}^{2}}$。

② 离子在匀强电场区域 BC 中做往返运动，设加速度为 $a$，则 $qE=ma$，$t_{2}=\frac{L}{v}+\frac{2v}{a}$，解得 $\frac{q}{m}=\frac{1}{2U}\left(L+\frac{4U}{E}\right)^{2}\frac{1}{t_{2}^{2}}$。

（2）两离子初速度分别为 $v$、$v'$，则 $t=\frac{L}{v}+\frac{2v}{\frac{qE}{m}}$，$t'=\frac{L}{v'}+\frac{2v'}{\frac{qE}{m}}$，$\Delta t=t-t'=\left(\frac{L}{vv'}-\frac{2m}{qE}\right)(v'-v)$。要使 $\Delta t=0$，则须 $\frac{L}{vv'}-\frac{2m}{qE}=0$，所以 $E=\frac{2mvv'}{qL}$。
}

\item
%% number: 25
%% paperName: 2007年重庆理综第 25 题 20 分
%% typeId: 6
%% type: 计算题
%% chapter: 运动和力
%% point: 动量和能量
%% method: 无
%% score: 20
%% degree: 250
%% duplicateId: 0
%% body:
某兴趣小组设计了一种实验装置，用来研究碰撞问题，其模型如图所示，用完全相同的轻绳将 $n$ 个大小相同、质量不等的小球并列悬挂于一水平杆，球间有微小间隔，从左到右，球的编号依次为 1、2、3、……、$n$，球的质量依次递减，每球质量与其相邻左球质量之比为 $k$（$k<1$）。将 1 号球向左拉起，然后由静止释放，使其与 2 号球碰撞，2 号球再与 3 号球碰撞……所有碰撞皆为无机械能损失的正碰。（不计空气阻力，忽略绳的伸长，$g$ 取 $10\Umsq$）

\begin{enumerate}
	\item
	设与 $n+1$ 号球碰撞前，$n$ 号球的速度为 $v_{n}$，求 $n+1$ 号球碰撞后的速度；
	\item
	若 $n=5$，在 1 号球向左拉高 $h$ 的情况下，要使 5 号球碰撞后升高 $16h$（$16h$ 小于绳长），问 $k$ 值为多少？
	\item
	在第（2）问的条件下，悬挂哪个球的绳最容易断，为什么？
\end{enumerate}

\onepicture[w=4.0cm, align=right]{figs/25.png}

%% answer:
\jdanswer{
\begin{enumerate}
	\item
	$\frac{2v_{n}}{k+1}$。

	\item
	$k=\sqrt{2}-1\approx0.414$。

	\item
	由题意 1 号球的重力最大，又由机械能守恒可知 1 号球在最低点碰前的动能也最大，张力最大，故悬挂 1 号球的绳容易断。
\end{enumerate}
}

%% memo:
\memoanswer{
（1）设 $n$ 号球质量为 $m_{n}$，$n+1$ 号球质量为 $m_{n+1}$，碰撞后的速度分别为 $v_{n}'$、$v_{n+1}'$，取水平向右为正方向。据题意，$n$ 号球与 $n+1$ 号球碰撞前的速度分别为 $v_{n}$、0，且 $m_{n+1}=km_{n}$。根据动量守恒有 $m_{n}v_{n}=m_{n}v_{n}'+km_{n}v_{n+1}'$，根据机械能守恒有 $\frac{1}{2}m_{n}v_{n}^{2}=\frac{1}{2}m_{n}v_{n}'^{2}+\frac{1}{2}km_{n}v_{n+1}'^{2}$，解得 $v_{n+1}'=\frac{2v_{n}}{k+1}$（$v_{n+1}'=0$ 舍去）。

（2）设 1 号球摆至最低点时的速度为 $v_{1}$，由机械能守恒定律有 $m_{1}gh=\frac{1}{2}m_{1}v_{1}^{2}$，得 $v_{1}=\sqrt{2gh}$。5 号球碰后瞬间的速度 $v_{5}=\left(\frac{2}{1+k}\right)^{4}v_{1}$，又 5 号球碰后升高 $16h$，则 $v_{5}=\sqrt{2g\cdot16h}=4\sqrt{2gh}$，故 $\left(\frac{2}{1+k}\right)^{4}=4$，解得 $k=\sqrt{2}-1\approx0.414$。

（3）设绳长为 $l$，每个球在最低点时，细绳对球的拉力为 $F$，由牛顿第二定律有 $F-m_{n}g=m_{n}\frac{v_{n}^{2}}{l}$，则 $F=m_{n}g+m_{n}\frac{v_{n}^{2}}{l}=m_{n}g+\frac{2}{l}E_{kn}$，式中 $E_{kn}$ 为 $n$ 号球在最低点的动能。由题意 1 号球的重力最大，又由机械能守恒可知 1 号球在最低点碰前的动能也最大，根据上式可判断在 1 号球碰前瞬间悬挂 1 号球细绳的张力最大，故悬挂 1 号球的绳容易断。
}

\end{enumerate}

\end{document}
